\documentclass[12pt]{article}
\usepackage[utf8]{inputenc}
\usepackage[T1]{fontenc}
\usepackage[margin=1in]{geometry}
\usepackage{setspace}
\usepackage{xcolor}
\usepackage[colorlinks=true, urlcolor=blue!60!black, linkcolor=blue!60!black]{hyperref}
\setlength{\parindent}{2em}
\setlength{\parskip}{1.25em}
\renewcommand{\baselinestretch}{1.0}
\begin{document}
\begin{center}
{\LARGE \textit{Balzac v. Porto Rico}, 258 U.S. 298 (1922)\par}
\vspace{0.8em}
{\large Case Brief --- Four-Page Edition\par}
\end{center}
\vspace{1em}

\section{Memory Jogger}

A Puerto Rican newspaper editor convicted of criminal libel without a jury argued that the 1917 grant of U.S. citizenship carried the Sixth Amendment with it; a unanimous Court held that Puerto Rico remained an unincorporated territory, where only fundamental rights apply of their own force and the jury right does not, because locality rather than citizenship controls (\href{https://supreme.justia.com/cases/federal/us/258/298/}{Justia}).

\section{Detailed Case Facts}

Jesus M. Balzac edited \textit{El Baluarte}, a daily newspaper in Arecibo, Porto Rico. The Jones Act of March 2, 1917, conferred U.S. citizenship on citizens of Porto Rico and enacted a statutory bill of rights that omitted grand and petit juries. On April 16 and 23, 1918, Balzac published two articles attacking the conduct of the island's Governor.

\section{Procedural History}

The district attorney filed two informations charging criminal libel in the District Court for Arecibo. Criminal libel was a misdemeanor, and local procedure granted juries only in felony cases (a 1919 statute extending juries to misdemeanors came too late). Balzac demanded a jury under the Sixth Amendment as an American citizen and claimed First Amendment protection for fair comment. The court overruled both defenses, tried him without a jury, and sentenced him to five months and four months in jail. The Supreme Court of Porto Rico affirmed, and Balzac brought writs of error. The U.S. Supreme Court affirmed on April 10, 1922.

\section{Judicial Votes}

Affirmed, \textbf{9--0}.
\begin{itemize}
\item \textbf{Majority (8):} Chief Justice Taft, joined by Justices McKenna, Day, Van Devanter, Pitney, McReynolds, Brandeis, and Clarke.
\item \textbf{Concurrence:} Justice Holmes, in the result, without opinion. No dissents.
\end{itemize}

\section{Holding}

The Jones Act did not incorporate Porto Rico into the United States. Under the \textit{Insular Cases} and the Territories Clause, only ``fundamental'' constitutional rights apply automatically in unincorporated territories; the Sixth Amendment jury right is not among them; and it is locality, not citizenship, that determines whether the jury guarantees apply. The Court also rejected the First Amendment defense, calling the articles ``excessive and outrageous.''

\section{Analysis of Opinions}

\textbf{Majority (Taft).} After confirming jurisdiction, Taft consolidated the \textit{Insular Cases}. \emph{Incorporation doctrine:} Justice White's \textit{Downes v. Bidwell} concurrence, confirmed in \textit{Dorr}, had become ``the settled law of the court'': jury guarantees apply in incorporated territories destined for statehood, not in those ``merely belonging to'' the United States. \emph{Clear statement:} Neither the Jones Act's title nor its text declared incorporation, and its statutory bill of rights deliberately omitted juries, showing that Congress did not think the federal Bill of Rights applied of its own force. \emph{Citizenship:} Congress granted citizenship to answer Puerto Ricans' ``yearning'' for diplomatic protection and to let them move to the mainland and enjoy full rights there; within the island, even a mainland citizen had no federal jury right. \emph{Alaska distinguished:} \textit{Rasmussen} treated Alaska as incorporated as a contiguous, sparsely settled territory, while Porto Rico was a ``distant ocean'' community of ``different origin and language'' trained in civil law, so Congress could let residents decide when to adopt juries. \emph{Courts:} The federal court in Porto Rico was a legislative territorial court, not a ``true United States court'' under Article~III. ``Certain fundamental personal rights'' such as due process always applied, but jury trial did not.

\textbf{Holmes (in the result).} Concurred without opinion; scholars read his silence as discomfort with Taft's rationale.

\textbf{Later developments.} \textit{Reid v. Covert} (1957) refused to extend the \textit{Insular Cases} to courts-martial of civilians; \textit{Flores de Otero} and \textit{Torres} applied equal protection, due process, and the Fourth Amendment in Puerto Rico; \textit{Harris v. Rosario} upheld lower benefits. In \textit{United States v. Vaello Madero} (2022), Justice Gorsuch urged overruling the \textit{Insular Cases} as relics rooted in racial stereotypes.

\section{Examples: Future Applications}

\begin{enumerate}
\item \textbf{Same-Side: Citizenship Without Grand Jury.} Congress grants citizenship in a new unincorporated Pacific territory with a statutory bill of rights omitting grand juries. Absent express incorporation, a felony prosecution without indictment stands.
\item \textbf{Same-Side: Mainland-Born Citizen in San Juan.} An Ohio-born citizen tried for a misdemeanor in 1920 Porto Rico demands a jury. Locality, not status, controls, so he may be tried without one, as Taft noted.
\item \textbf{Opposite-Side: Incorporation or a Fundamental Right.} Congress expressly incorporates a territory as a future State, or a territory authorizes suspicionless home searches. The full Constitution applies in the first case, as in \textit{Rasmussen}; fundamental rights constrain the second, as in \textit{Torres}.
\item \textbf{Opposite-Side: Court-Martial of a Civilian Abroad.} A soldier's civilian wife is court-martialed overseas for killing her husband. \textit{Reid v. Covert} refused to extend the \textit{Insular Cases}.
\item \textbf{Fence-Sitter: Felony Bench Trial Today.} A territorial court tries a serious felony without a jury. \textit{Balzac} has never been overruled, but \textit{Duncan v. Louisiana}, \textit{Reid}, and Justice Gorsuch's \textit{Vaello Madero} concurrence make the Supreme Court's answer uncertain.
\end{enumerate}

\section{Critique}

\textbf{Text and structure.} The Territories Clause grants power to make ``needful Rules and Regulations,'' not to suspend the Bill of Rights, and, as Justice Gorsuch noted, the Constitution nowhere distinguishes ``incorporated'' from ``unincorporated'' territories.

\textbf{Doctrinal stranding.} Taft's premise that jury trial is not fundamental collides with \textit{Duncan v. Louisiana} (1968), which held the jury right ``fundamental to the American scheme of justice'' and incorporated it against the States, leaving \textit{Balzac}'s core premise stranded and its cultural assumptions exposed.

\section{Key Quotations}

\begin{quote}
``It is locality that is determinative of the application of the Constitution, in such matters as judicial procedure, and not the status of the people who live in it.''
\end{quote}
\noindent --- Chief Justice Taft, 258 U.S. at 309.

\begin{quote}
``Before that, the purpose of Congress might well be a matter of mere inference from various legislative acts; but in these latter days, incorporation is not to be assumed without express declaration, or an implication so strong as to exclude any other view.''
\end{quote}
\noindent --- Chief Justice Taft, 258 U.S. at 306.

\begin{quote}
``Congress has thought that a people like the Filipinos, or the Porto Ricans, trained to a complete judicial system which knows no juries, living in compact and ancient communities, with definitely formed customs and political conceptions, should be permitted themselves to determine how far they wish to adopt this institution of Anglo-Saxon origin, and when.''
\end{quote}
\noindent --- Chief Justice Taft, 258 U.S. at 310.

\end{document}
